\documentclass[11pt,a4paper]{article}
\usepackage[utf8]{inputenc}
\usepackage[T1]{fontenc}
\usepackage[catalan]{babel}
\usepackage{lmodern}
\usepackage[sfdefault,lining]{FiraSans}
\usepackage[a4paper,top=2.3cm,bottom=2.2cm,left=1.8cm,right=1.8cm,headheight=1.3cm,headsep=0.4cm,footskip=1cm]{geometry}
\usepackage{amsmath,amssymb}
\usepackage{icomma}
\usepackage{graphicx}
\usepackage[table]{xcolor}
\usepackage{tikz}
\usetikzlibrary{calc}
\usepackage{fontawesome5}
\usepackage{ccicons}
\usepackage{fancyhdr}
\usepackage{lastpage}
\usepackage{tasks}
\usepackage{needspace}
\usepackage[most]{tcolorbox}
\usepackage{array}
\usepackage[hidelinks,pdftitle={Fitxa d'exercicis de la unitat 3: Programació lineal. 2n de Batxillerat CCSS},pdfsubject={Matemàtiques Aplicades a les Ciències Socials II},pdfkeywords={programació lineal, inequacions, sistemes d'inequacions, regió factible, vèrtexs, funció objectiu, rectes de nivell, PAU}]{hyperref}

\definecolor{navy}{HTML}{1B2748}
\definecolor{accent}{HTML}{E8603A}
\definecolor{teal}{HTML}{1F9E90}
\definecolor{soft}{HTML}{EEF1F8}
\definecolor{muted}{HTML}{5D6685}

\setlength{\parindent}{0pt}
\setlength{\parskip}{3pt}
\settasks{label=\alph*), label-format=\bfseries\color{navy}, label-width=1.3em, item-indent=1.8em, column-sep=0.6em, after-item-skip=5pt, before-skip=3pt, after-skip=4pt}

% ---------- capçalera i peu ----------
\newcommand{\llicencia}{\href{https://creativecommons.org/licenses/by-nc-sa/4.0/deed.ca}{\ccbyncsa}\ \ Llicència Creative Commons Reconeixement-NoComercial-CompartirIgual 4.0 Internacional (CC BY-NC-SA 4.0)}
\pagestyle{fancy}
\fancyhf{}
\fancyhead[L]{\raisebox{-0.35\height}{\includegraphics[height=0.95cm]{logo-gva.png}}}
\fancyhead[C]{\small\color{muted}Fitxa d'exercicis · Unitat 3: Programació lineal · 2n Batxillerat CCSS}
\fancyhead[R]{\raisebox{-0.35\height}{\includegraphics[height=0.95cm]{logo-vermellar.png}}}
\fancyfoot[L]{\scriptsize\color{muted}\llicencia}
\fancyfoot[R]{\small\color{muted}\thepage}
\renewcommand{\headrulewidth}{0.4pt}
\renewcommand{\headrule}{\color{muted!40}\hrule height \headrulewidth}
\fancypagestyle{primera}{\fancyhf{}\renewcommand{\headrulewidth}{0pt}
  \fancyfoot[L]{\scriptsize\color{muted}\llicencia}
  \fancyfoot[R]{\small\color{muted}\thepage}}

% ---------- nivells i exercicis ----------
\newcommand{\estrelles}[1]{{\color{accent}\ifcase#1\or\faStar\faStar[regular]\faStar[regular]\or\faStar\faStar\faStar[regular]\or\faStar\faStar\faStar\fi}}
\newcounter{exer}
\newcommand{\ex}[1]{\par\needspace{4\baselineskip}\medskip\refstepcounter{exer}\noindent{\color{navy}\bfseries\theexer.}\ \ }

\newcommand{\seccio}[2]{\par\needspace{10\baselineskip}\bigskip
  \begin{tcolorbox}[colback=navy,colframe=navy,coltext=white,arc=2mm,boxrule=0pt,left=3mm,right=3mm,top=1.2mm,bottom=1.2mm]
  {\large\bfseries #1\quad #2}
  \end{tcolorbox}\vspace{-2pt}}


% quadre de recordatori (només a la fitxa d'exercicis)
\newcommand{\recorda}[1]{\par\smallskip\begin{tcolorbox}[colback=accent!7,colframe=accent!55,arc=2mm,boxrule=0.5pt,left=2.5mm,right=2.5mm,top=1mm,bottom=1mm]\small{\color{accent}\faLightbulb[regular]}\ \textbf{Recorda:}\ #1\end{tcolorbox}}
% etiqueta dels exercicis de la PAU
\newcommand{\pau}[1]{{\color{teal}\small\bfseries\faUniversity\ PAU · #1.}\ }
\newcommand{\res}[1]{#1}\newcommand{\rest}[1]{#1}

% =====================================================================
% Capçalera FentMates de les fitxes: un color per a cada curs.
% Els colors són els de la pàgina inicial (LEVELS a assets/js/materials.js);
% el logo ve de _eines/logo/pdf/ (python3 _eines/logo/genera.py).
%
% Ús, al final del preàmbul de la unitat:
%   % =====================================================================
% Capçalera FentMates de les fitxes: un color per a cada curs.
% Els colors són els de la pàgina inicial (LEVELS a assets/js/materials.js);
% el logo ve de _eines/logo/pdf/ (python3 _eines/logo/genera.py).
%
% Ús, al final del preàmbul de la unitat:
%   % =====================================================================
% Capçalera FentMates de les fitxes: un color per a cada curs.
% Els colors són els de la pàgina inicial (LEVELS a assets/js/materials.js);
% el logo ve de _eines/logo/pdf/ (python3 _eines/logo/genera.py).
%
% Ús, al final del preàmbul de la unitat:
%   \input{../comu/fentmates.tex}
%   \fmcurs{1eso}
% i, al començament del document:
%   \capcalera{Matemàtiques}{1r ESO}{Fitxa d'exercicis · Unitat 2}{Divisibilitat}
% Opcional abans de \input: \def\fmneed{14} (espai mínim abans de cada secció).
% =====================================================================
\makeatletter
\def\fm@logodir{../../logo/pdf}
\providecommand{\fmneed}{10}
\def\fm@cursnom{}

% clau del curs, color i color del text damunt del color
\newcommand{\fm@curs}[3]{\@namedef{fm@color@#1}{#2}\@namedef{fm@text@#1}{#3}}
\fm@curs{1eso}{E6B84A}{navy}
\fm@curs{2eso}{3B82F6}{white}
\fm@curs{3eso}{FF7A50}{white}
\fm@curs{4eso-a}{A78BFA}{white}
\fm@curs{4eso-b}{6D28D9}{white}
\fm@curs{1batx-mi}{9B1C1C}{white}
\fm@curs{1batx-ccss}{EF6B6B}{white}
\fm@curs{2batx-mii}{1E6B47}{white}
\fm@curs{2batx-ccss}{2BB3A3}{white}

\newcommand{\fmcurs}[1]{%
  \@ifundefined{fm@color@#1}{\PackageError{fentmates}{Curs desconegut: #1}{}}{}%
  \edef\fm@tmp{\noexpand\definecolor{curs}{HTML}{\@nameuse{fm@color@#1}}%
               \noexpand\colorlet{curstext}{\@nameuse{fm@text@#1}}}\fm@tmp
  \colorlet{cursfosc}{curs!72!black}%
  \def\fm@logo{\fm@logodir/vertical-#1.pdf}%
  \def\fm@marca{\fm@logodir/marca-#1.pdf}}

% ---------- capçalera de la primera pàgina ----------
% #1 matèria (petit), #2 curs, #3 tipus de document i unitat, #4 títol
\newcommand{\capcalera}[4]{%
  \gdef\fm@cursnom{#2}%
  \par\noindent
  \begin{tcolorbox}[colback=curs,colframe=curs,boxrule=0pt,arc=2.2mm,width=3.5cm,height=2.15cm,
      valign=center,left=3mm,right=2mm,top=1mm,bottom=1mm,nobeforeafter,box align=center]
    \color{curstext}{\footnotesize #1}\par\vspace{1pt}{\fontsize{19}{21}\selectfont\bfseries #2\par}%
  \end{tcolorbox}\hspace{5mm}%
  \begin{minipage}[c]{\dimexpr\linewidth-3.5cm-5mm-2.9cm\relax}
    {\color{cursfosc}\small\bfseries\MakeUppercase{#3}}\par\vspace{2pt}
    {\color{navy}\fontsize{20}{23}\selectfont\bfseries #4\par}
  \end{minipage}\hfill
  \begin{minipage}[c]{2.6cm}\raggedleft\includegraphics[width=2.6cm]{\fm@logo}\end{minipage}\par
  \vspace{2pt}}

% ---------- barres de secció del color del curs ----------
\renewcommand{\seccio}[2]{\par\needspace{\fmneed\baselineskip}\bigskip
  \begin{tcolorbox}[colback=curs,colframe=curs,coltext=curstext,arc=2mm,boxrule=0pt,left=3mm,right=3mm,top=1.2mm,bottom=1.2mm]
  {\large\bfseries #1\quad #2}
  \end{tcolorbox}\vspace{-2pt}}

% ---------- capçalera i peu de les pàgines ----------
% primera pàgina: logos institucionals al peu
\fancypagestyle{primera}{\fancyhf{}\renewcommand{\headrulewidth}{0pt}
  \fancyfoot[L]{\raisebox{-0.32\height}{\includegraphics[height=0.78cm]{logo-gva.png}}\hspace{5mm}%
    \raisebox{-0.32\height}{\includegraphics[height=0.78cm]{logo-vermellar.png}}}
  \fancyfoot[R]{\scriptsize\color{muted}\href{https://creativecommons.org/licenses/by-nc-sa/4.0/deed.ca}{\ccbyncsa}\ \ CC BY-NC-SA 4.0 · fentmates.org\qquad\small\thepage}}
% la resta: el símbol a l'esquerra i el curs a la dreta, amb el títol al centre com abans
\fancyhead[L]{\raisebox{-0.3\height}{\includegraphics[height=0.62cm]{\fm@marca}}}
\fancyhead[R]{\small\bfseries\color{cursfosc}\fm@cursnom}
\renewcommand{\headrule}{\color{curs!45}\hrule height \headrulewidth}
\makeatother

%   \fmcurs{1eso}
% i, al començament del document:
%   \capcalera{Matemàtiques}{1r ESO}{Fitxa d'exercicis · Unitat 2}{Divisibilitat}
% Opcional abans de \input: \def\fmneed{14} (espai mínim abans de cada secció).
% =====================================================================
\makeatletter
\def\fm@logodir{../../logo/pdf}
\providecommand{\fmneed}{10}
\def\fm@cursnom{}

% clau del curs, color i color del text damunt del color
\newcommand{\fm@curs}[3]{\@namedef{fm@color@#1}{#2}\@namedef{fm@text@#1}{#3}}
\fm@curs{1eso}{E6B84A}{navy}
\fm@curs{2eso}{3B82F6}{white}
\fm@curs{3eso}{FF7A50}{white}
\fm@curs{4eso-a}{A78BFA}{white}
\fm@curs{4eso-b}{6D28D9}{white}
\fm@curs{1batx-mi}{9B1C1C}{white}
\fm@curs{1batx-ccss}{EF6B6B}{white}
\fm@curs{2batx-mii}{1E6B47}{white}
\fm@curs{2batx-ccss}{2BB3A3}{white}

\newcommand{\fmcurs}[1]{%
  \@ifundefined{fm@color@#1}{\PackageError{fentmates}{Curs desconegut: #1}{}}{}%
  \edef\fm@tmp{\noexpand\definecolor{curs}{HTML}{\@nameuse{fm@color@#1}}%
               \noexpand\colorlet{curstext}{\@nameuse{fm@text@#1}}}\fm@tmp
  \colorlet{cursfosc}{curs!72!black}%
  \def\fm@logo{\fm@logodir/vertical-#1.pdf}%
  \def\fm@marca{\fm@logodir/marca-#1.pdf}}

% ---------- capçalera de la primera pàgina ----------
% #1 matèria (petit), #2 curs, #3 tipus de document i unitat, #4 títol
\newcommand{\capcalera}[4]{%
  \gdef\fm@cursnom{#2}%
  \par\noindent
  \begin{tcolorbox}[colback=curs,colframe=curs,boxrule=0pt,arc=2.2mm,width=3.5cm,height=2.15cm,
      valign=center,left=3mm,right=2mm,top=1mm,bottom=1mm,nobeforeafter,box align=center]
    \color{curstext}{\footnotesize #1}\par\vspace{1pt}{\fontsize{19}{21}\selectfont\bfseries #2\par}%
  \end{tcolorbox}\hspace{5mm}%
  \begin{minipage}[c]{\dimexpr\linewidth-3.5cm-5mm-2.9cm\relax}
    {\color{cursfosc}\small\bfseries\MakeUppercase{#3}}\par\vspace{2pt}
    {\color{navy}\fontsize{20}{23}\selectfont\bfseries #4\par}
  \end{minipage}\hfill
  \begin{minipage}[c]{2.6cm}\raggedleft\includegraphics[width=2.6cm]{\fm@logo}\end{minipage}\par
  \vspace{2pt}}

% ---------- barres de secció del color del curs ----------
\renewcommand{\seccio}[2]{\par\needspace{\fmneed\baselineskip}\bigskip
  \begin{tcolorbox}[colback=curs,colframe=curs,coltext=curstext,arc=2mm,boxrule=0pt,left=3mm,right=3mm,top=1.2mm,bottom=1.2mm]
  {\large\bfseries #1\quad #2}
  \end{tcolorbox}\vspace{-2pt}}

% ---------- capçalera i peu de les pàgines ----------
% primera pàgina: logos institucionals al peu
\fancypagestyle{primera}{\fancyhf{}\renewcommand{\headrulewidth}{0pt}
  \fancyfoot[L]{\raisebox{-0.32\height}{\includegraphics[height=0.78cm]{logo-gva.png}}\hspace{5mm}%
    \raisebox{-0.32\height}{\includegraphics[height=0.78cm]{logo-vermellar.png}}}
  \fancyfoot[R]{\scriptsize\color{muted}\href{https://creativecommons.org/licenses/by-nc-sa/4.0/deed.ca}{\ccbyncsa}\ \ CC BY-NC-SA 4.0 · fentmates.org\qquad\small\thepage}}
% la resta: el símbol a l'esquerra i el curs a la dreta, amb el títol al centre com abans
\fancyhead[L]{\raisebox{-0.3\height}{\includegraphics[height=0.62cm]{\fm@marca}}}
\fancyhead[R]{\small\bfseries\color{cursfosc}\fm@cursnom}
\renewcommand{\headrule}{\color{curs!45}\hrule height \headrulewidth}
\makeatother

%   \fmcurs{1eso}
% i, al començament del document:
%   \capcalera{Matemàtiques}{1r ESO}{Fitxa d'exercicis · Unitat 2}{Divisibilitat}
% Opcional abans de \input: \def\fmneed{14} (espai mínim abans de cada secció).
% =====================================================================
\makeatletter
\def\fm@logodir{../../logo/pdf}
\providecommand{\fmneed}{10}
\def\fm@cursnom{}

% clau del curs, color i color del text damunt del color
\newcommand{\fm@curs}[3]{\@namedef{fm@color@#1}{#2}\@namedef{fm@text@#1}{#3}}
\fm@curs{1eso}{E6B84A}{navy}
\fm@curs{2eso}{3B82F6}{white}
\fm@curs{3eso}{FF7A50}{white}
\fm@curs{4eso-a}{A78BFA}{white}
\fm@curs{4eso-b}{6D28D9}{white}
\fm@curs{1batx-mi}{9B1C1C}{white}
\fm@curs{1batx-ccss}{EF6B6B}{white}
\fm@curs{2batx-mii}{1E6B47}{white}
\fm@curs{2batx-ccss}{2BB3A3}{white}

\newcommand{\fmcurs}[1]{%
  \@ifundefined{fm@color@#1}{\PackageError{fentmates}{Curs desconegut: #1}{}}{}%
  \edef\fm@tmp{\noexpand\definecolor{curs}{HTML}{\@nameuse{fm@color@#1}}%
               \noexpand\colorlet{curstext}{\@nameuse{fm@text@#1}}}\fm@tmp
  \colorlet{cursfosc}{curs!72!black}%
  \def\fm@logo{\fm@logodir/vertical-#1.pdf}%
  \def\fm@marca{\fm@logodir/marca-#1.pdf}}

% ---------- capçalera de la primera pàgina ----------
% #1 matèria (petit), #2 curs, #3 tipus de document i unitat, #4 títol
\newcommand{\capcalera}[4]{%
  \gdef\fm@cursnom{#2}%
  \par\noindent
  \begin{tcolorbox}[colback=curs,colframe=curs,boxrule=0pt,arc=2.2mm,width=3.5cm,height=2.15cm,
      valign=center,left=3mm,right=2mm,top=1mm,bottom=1mm,nobeforeafter,box align=center]
    \color{curstext}{\footnotesize #1}\par\vspace{1pt}{\fontsize{19}{21}\selectfont\bfseries #2\par}%
  \end{tcolorbox}\hspace{5mm}%
  \begin{minipage}[c]{\dimexpr\linewidth-3.5cm-5mm-2.9cm\relax}
    {\color{cursfosc}\small\bfseries\MakeUppercase{#3}}\par\vspace{2pt}
    {\color{navy}\fontsize{20}{23}\selectfont\bfseries #4\par}
  \end{minipage}\hfill
  \begin{minipage}[c]{2.6cm}\raggedleft\includegraphics[width=2.6cm]{\fm@logo}\end{minipage}\par
  \vspace{2pt}}

% ---------- barres de secció del color del curs ----------
\renewcommand{\seccio}[2]{\par\needspace{\fmneed\baselineskip}\bigskip
  \begin{tcolorbox}[colback=curs,colframe=curs,coltext=curstext,arc=2mm,boxrule=0pt,left=3mm,right=3mm,top=1.2mm,bottom=1.2mm]
  {\large\bfseries #1\quad #2}
  \end{tcolorbox}\vspace{-2pt}}

% ---------- capçalera i peu de les pàgines ----------
% primera pàgina: logos institucionals al peu
\fancypagestyle{primera}{\fancyhf{}\renewcommand{\headrulewidth}{0pt}
  \fancyfoot[L]{\raisebox{-0.32\height}{\includegraphics[height=0.78cm]{logo-gva.png}}\hspace{5mm}%
    \raisebox{-0.32\height}{\includegraphics[height=0.78cm]{logo-vermellar.png}}}
  \fancyfoot[R]{\scriptsize\color{muted}\href{https://creativecommons.org/licenses/by-nc-sa/4.0/deed.ca}{\ccbyncsa}\ \ CC BY-NC-SA 4.0 · fentmates.org\qquad\small\thepage}}
% la resta: el símbol a l'esquerra i el curs a la dreta, amb el títol al centre com abans
\fancyhead[L]{\raisebox{-0.3\height}{\includegraphics[height=0.62cm]{\fm@marca}}}
\fancyhead[R]{\small\bfseries\color{cursfosc}\fm@cursnom}
\renewcommand{\headrule}{\color{curs!45}\hrule height \headrulewidth}
\makeatother

\fmcurs{2batx-ccss}

\begin{document}
\thispagestyle{primera}

\capcalera{Matemàtiques CCSS}{2n Batx}{Fitxa d'exercicis · Unitat 3}{Programació lineal}
\vspace{10pt}
Nom i cognoms: \rule{8.4cm}{0.4pt}\hfill Grup: \rule{1.6cm}{0.4pt}\hfill Data: \rule{2.2cm}{0.4pt}
\vspace{8pt}

\begin{tcolorbox}[colback=soft,colframe=soft,arc=2mm,left=3mm,right=3mm,top=1.5mm,bottom=1.5mm]
\small Fes les gràfiques amb regle i paper quadriculat, i indica sempre els vèrtexs de la regió. Els exercicis marcats amb {\color{teal}\bfseries\faUniversity\ PAU} són de les proves d'accés a la universitat de la Comunitat Valenciana; en trobaràs molts més, resolts, a la pàgina d'exercicis PAU de la unitat.
\end{tcolorbox}

% =====================================================
\seccio{1}{Inequacions lineals amb dues incògnites}

\ex{1} Esbrina quins dels punts següents són solució de la inequació $3x-2y\leq4$:\begin{tasks}(4)
  \task $(0,\,0)$
  \task $(2,\,1)$
  \task $(-1,\,3)$
  \task $(4,\,-2)$
\end{tasks}\recorda{Un punt $(x_0,y_0)$ és solució si, en substituir-lo, la desigualtat és certa.}

\ex{1} Resol gràficament les inequacions següents:\begin{tasks}(4)
  \task $x+y\leq4$
  \task $2x-y>2$
  \task $y\geq3x$
  \task $x-3y<6$
\end{tasks}\recorda{1) Dibuixa la recta frontera: contínua si és $\leq,\geq$ i discontínua si és $<,>$. \ 2) Tria un punt de prova que no siga de la recta (normalment $(0,\,0)$). \ 3) Si el compleix, la solució és el semipla que el conté; si no, l'altre.}

\ex{1} Resol gràficament les inequacions següents, en les quals només apareix una de les incògnites:\begin{tasks}(4)
  \task $x\geq-2$
  \task $y<3$
  \task $2x\leq5$
  \task $-y\geq1$
\end{tasks}\recorda{La recta $x=k$ és \textbf{vertical} i $y=k$ és \textbf{horitzontal}. \ $x\geq k$: a la dreta; \ $x\leq k$: a l'esquerra; \ $y\geq k$: per damunt; \ $y\leq k$: per davall. No cal punt de prova.}

\ex{2} Escriu la inequació que té com a solució cada semipla ombrejat:\begin{tasks}(2)
  \task \begin{tikzpicture}[x=0.5250cm,y=0.5250cm,baseline=(current bounding box.center)]\begin{scope}\clip (-2,-2) rectangle (6,4);\foreach \i in {-2,0,...,6} \draw[muted!25,very thin] (\i,-2)--(\i,4);\foreach \i in {-2,-1,...,4} \draw[muted!25,very thin] (-2,\i)--(6,\i);\fill[teal!22] (-2.0000,-2.0000)--(6.0000,-2.0000)--(6.0000,-1.0000)--(-2.0000,3.0000)--cycle;\end{scope}\draw[-latex,thick] (-2,0)--(6.5,0) node[below left=-1pt,font=\scriptsize]{$x$};\draw[-latex,thick] (0,-2)--(0,4.25) node[below left=-1pt,font=\scriptsize]{$y$};\node[below,font=\tiny,text=muted] at (-2,0) {$-2$};\node[below,font=\tiny,text=muted] at (2,0) {$2$};\node[below,font=\tiny,text=muted] at (4,0) {$4$};\node[below,font=\tiny,text=muted] at (6,0) {$6$};\node[left,font=\tiny,text=muted] at (0,-2) {$-2$};\node[left,font=\tiny,text=muted] at (0,-1) {$-1$};\node[left,font=\tiny,text=muted] at (0,1) {$1$};\node[left,font=\tiny,text=muted] at (0,2) {$2$};\node[left,font=\tiny,text=muted] at (0,3) {$3$};\node[left,font=\tiny,text=muted] at (0,4) {$4$};\begin{scope}\clip (-2,-2) rectangle (6,4);\draw[accent,line width=1.1pt] (-3.0000,3.5000)--(7.0000,-1.5000);\end{scope}\end{tikzpicture}
  \task \begin{tikzpicture}[x=0.7000cm,y=0.5750cm,baseline=(current bounding box.center)]\begin{scope}\clip (-2,-3) rectangle (4,5);\foreach \i in {-2,-1,...,4} \draw[muted!25,very thin] (\i,-3)--(\i,5);\foreach \i in {-2,0,...,5} \draw[muted!25,very thin] (-2,\i)--(4,\i);\fill[teal!22] (-2.0000,-3.0000)--(-1.0000,-3.0000)--(3.0000,5.0000)--(-2.0000,5.0000)--cycle;\end{scope}\draw[-latex,thick] (-2,0)--(4.25,0) node[below left=-1pt,font=\scriptsize]{$x$};\draw[-latex,thick] (0,-3)--(0,5.5) node[below left=-1pt,font=\scriptsize]{$y$};\node[below,font=\tiny,text=muted] at (-2,0) {$-2$};\node[below,font=\tiny,text=muted] at (-1,0) {$-1$};\node[below,font=\tiny,text=muted] at (1,0) {$1$};\node[below,font=\tiny,text=muted] at (2,0) {$2$};\node[below,font=\tiny,text=muted] at (3,0) {$3$};\node[below,font=\tiny,text=muted] at (4,0) {$4$};\node[left,font=\tiny,text=muted] at (0,-2) {$-2$};\node[left,font=\tiny,text=muted] at (0,2) {$2$};\node[left,font=\tiny,text=muted] at (0,4) {$4$};\begin{scope}\clip (-2,-3) rectangle (4,5);\draw[accent,line width=1.1pt,dashed] (-1.5000,-4.0000)--(3.5000,6.0000);\end{scope}\end{tikzpicture}
\end{tasks}\recorda{Troba l'equació de la recta a partir de dos punts i decideix el signe amb un punt de l'ombrejat.}

\ex{2} Expressa cada frase amb una inequació i representa-la gràficament ($x$ i $y$ són dos nombres reals):\begin{tasks}(1)
  \task El doble d'un nombre més el triple d'un altre no supera 12.
  \task La diferència entre el primer nombre i el segon és, com a mínim, 2.
\end{tasks}


% =====================================================
\seccio{2}{Sistemes d'inequacions. Regió factible i vèrtexs}

\ex{1} Resol gràficament els sistemes següents i calcula els vèrtexs de la regió solució. Indica si és acotada o no:\begin{tasks}(2)
  \task $\left\{\begin{aligned}x+y&\leq5\\x-y&\geq-1\\x\geq0,\ y&\geq0\end{aligned}\right.$
  \task $\left\{\begin{aligned}2x+y&\geq4\\x+2y&\geq5\\x\geq0,\ y&\geq0\end{aligned}\right.$
\end{tasks}\recorda{La solució és la zona comuna a tots els semiplans. Els vèrtexs són els punts de tall de dues rectes frontera que compleixen \textbf{totes} les inequacions.}

\ex{1} Representa la regió definida per les inequacions següents i calcula'n els vèrtexs:\[1\leq x\leq4,\qquad 0\leq y\leq3,\qquad x+y\leq6\]\recorda{Una doble desigualtat $a\leq x\leq b$ són dues inequacions: $x\geq a$ i $x\leq b$ (una franja vertical).}

\ex{2} Resol gràficament els sistemes següents:\begin{tasks}(2)
  \task $\left\{\begin{aligned}x+y&\geq2\\x+y&\leq1\\x\geq0,\ y&\geq0\end{aligned}\right.$
  \task $\left\{\begin{aligned}y&\geq x\\y&\leq x+2\\0\leq x&\leq3\end{aligned}\right.$
\end{tasks}

\ex{2} Escriu el sistema d'inequacions que té com a solució la regió ombrejada:

\begin{center}\begin{tikzpicture}[x=0.6250cm,y=0.6250cm,baseline=(current bounding box.center)]\begin{scope}\clip (-1,-1) rectangle (7,5.5);\foreach \i in {0,2,...,7} \draw[muted!25,very thin] (\i,-1)--(\i,5.5);\foreach \i in {-1,0,...,5.5} \draw[muted!25,very thin] (-1,\i)--(7,\i);\fill[teal!22] (6.0000,0.0000)--(2.0000,4.0000)--(0.0000,0.0000)--cycle;\end{scope}\draw[-latex,thick] (-1,0)--(7.5,0) node[below left=-1pt,font=\scriptsize]{$x$};\draw[-latex,thick] (0,-1)--(0,5.75) node[below left=-1pt,font=\scriptsize]{$y$};\node[below,font=\tiny,text=muted] at (2,0) {$2$};\node[below,font=\tiny,text=muted] at (4,0) {$4$};\node[below,font=\tiny,text=muted] at (6,0) {$6$};\node[left,font=\tiny,text=muted] at (0,-1) {$-1$};\node[left,font=\tiny,text=muted] at (0,1) {$1$};\node[left,font=\tiny,text=muted] at (0,2) {$2$};\node[left,font=\tiny,text=muted] at (0,3) {$3$};\node[left,font=\tiny,text=muted] at (0,4) {$4$};\node[left,font=\tiny,text=muted] at (0,5) {$5$};\begin{scope}\clip (-1,-1) rectangle (7,5.5);\draw[accent,line width=1.1pt] (-1.0000,-2.0000)--(3.2500,6.5000);\draw[blue!70!black,line width=1.1pt] (8.0000,-2.0000)--(-0.5000,6.5000);\end{scope}\filldraw[fill=white,draw=navy] (6.0000,0.0000) circle (1.5pt);\node[font=\tiny,inner sep=1.5pt,fill=white,fill opacity=.75,text opacity=1] at ($(6.0000,0.0000)+(-21.8:0.36cm)$) {$(6,\,0)$};\filldraw[fill=white,draw=navy] (2.0000,4.0000) circle (1.5pt);\node[font=\tiny,inner sep=1.5pt,fill=white,fill opacity=.75,text opacity=1] at ($(2.0000,4.0000)+(104.0:0.36cm)$) {$(2,\,4)$};\filldraw[fill=white,draw=navy] (0.0000,0.0000) circle (1.5pt);\node[font=\tiny,inner sep=1.5pt,fill=white,fill opacity=.75,text opacity=1] at ($(0.0000,0.0000)+(-153.4:0.36cm)$) {$(0,\,0)$};\end{tikzpicture}\end{center}

\ex{2} Considera el sistema $\ x+2y\leq8,\ \ 3x+y\leq9,\ \ x\geq0,\ \ y\geq0$.\begin{tasks}(1)
  \task Esbrina si els punts $(1,\,2)$, $(2,\,3)$ i $(3,\,1)$ pertanyen a la regió factible.
  \task Representa la regió i calcula'n els vèrtexs.
\end{tasks}

\ex{2} Representa la regió definida per $\ x\geq1,\ \ y\geq2,\ \ x+y\geq5\ $. És acotada? Calcula'n els vèrtexs.


% =====================================================
\seccio{3}{Optimització d'una funció objectiu sobre una regió}

\ex{1} Calcula el màxim i el mínim de la funció $f(x,y)=3x+2y$ en la regió de l'apartat a) de l'exercici 6.\recorda{Si la regió és acotada, el màxim i el mínim s'assoleixen en vèrtexs: avalua la funció en tots.}

\ex{1} La regió factible d'un problema és el polígon de vèrtexs $A(0,\,0)$, $B(6,\,0)$, $C(5,\,3)$, $D(2,\,5)$ i $E(0,\,4)$. Troba el màxim i el mínim de les funcions $f(x,y)=x+2y$ i $g(x,y)=2x-y$ en aquesta regió.

\ex{2} Calcula, si existeixen, el màxim i el mínim de $f(x,y)=4x+3y$ en la regió de l'apartat b) de l'exercici 6.\recorda{Si la regió no és acotada, pot ser que no hi haja màxim (o mínim). Comprova-ho amb les rectes de nivell $4x+3y=k$.}

\ex{2} Resol el problema anterior gràficament amb rectes de nivell: dibuixa la recta $4x+3y=0$ i desplaça-la paral·lelament. Quin és el primer vèrtex de la regió que toca?


% =====================================================
\seccio{4}{Problemes de programació lineal}

\ex{1} Una fusteria fabrica taules i cadires. Cada taula necessita 3 hores de muntatge i 1 hora d'acabat, i cada cadira, 1 hora de muntatge i 1 hora d'acabat. Cada setmana es disposa de 60 hores de muntatge i 30 hores d'acabat. El benefici és de 40 € per taula i de 20 € per cadira. Quantes taules i quantes cadires cal fabricar per obtenir el màxim benefici?\recorda{Passos: variables (taula de dades) $\to$ funció objectiu $\to$ restriccions $\to$ regió factible $\to$ vèrtexs $\to$ avaluació $\to$ resposta.}

\ex{2} Un ramader vol barrejar dos pinsos, P i Q. Cada kg de P aporta 2 unitats de vitamina A i 1 de vitamina B, i cada kg de Q, 1 unitat de A i 3 de B. Cada dia el bestiar necessita almenys 8 unitats de A i 9 de B. Si el kg de P costa 3 € i el de Q, 2 €, quants kg de cada pinso ha de comprar perquè el cost siga mínim?

\ex{2} Una persona disposa de 30\,000 € per invertir en dos fons, A i B, amb rendibilitats anuals del 4\,\% i del 6\,\%. Vol posar almenys 10\,000 € en A i no vol invertir en B (més arriscat) més del que inverteix en A. Com ha de repartir els diners per obtenir el màxim rendiment?

\ex{3} Una empresa té dos magatzems: M1, amb 30 tones de producte, i M2, amb 25 tones. Ha de servir 20 tones a la botiga T1 i 25 a la botiga T2. El cost de transport per tona (en euros) és:

\begin{center}\small\begin{tabular}{c|cc}& T1 & T2\\\hline M1 & 5 & 3\\ M2 & 2 & 4\end{tabular}\end{center}Com s'ha d'organitzar el transport perquè el cost siga mínim? Quin és aquest cost?

\ex{2} \pau{C. Valenciana, juliol 2025} Una agència de viatges organitza excursions a la muntanya i a la platja. L'agència obté 700 euros de benefici per cada excursió a la muntanya i 500 euros per cada excursió a la platja. Disposa d'un total de 10 autobusos i 8 guies turístics. Cada excursió a la muntanya requereix 2 autobusos i 2 guies, mentre que cada excursió a la platja requereix 2 autobusos i 1 guia.\begin{tasks}(1)
  \task Quantes excursions a la muntanya i quantes a la platja ha d'organitzar per obtenir el màxim benefici possible?
  \task Quin és aquest benefici màxim?
\end{tasks}

\ex{2} \pau{C. Valenciana, juny 2025} Un creador de contingut desitja invertir en publicitat fins a 10\,000 euros. Ha considerat dues empreses, A i B. L'empresa A garanteix un increment anual en nombre de seguidors del 10\,\% del que s'hi inverteix, i la B, del 7\,\%. El creador vol invertir almenys 2000 euros en l'empresa B i un màxim de 6000 euros en la A, invertint sempre en la A almenys el que inverteix en la B.\begin{tasks}(1)
  \task Quants diners ha d'assignar a cada empresa perquè l'increment de seguidors siga màxim?
  \task Quin és aquest increment màxim?
\end{tasks}

\ex{2} \pau{C. Valenciana, juny 2023} El veterinari m'ha recomanat que el meu gos prenga diàriament un mínim de 8 unitats d'hidrats de carboni, 46 de proteïnes i 12 de greixos. Una llanda de la marca A conté 4 unitats d'hidrats, 6 de proteïnes i 1 de greixos; una llanda de la marca B, 2 unitats d'hidrats, 20 de proteïnes i 12 de greixos. La llanda A costa 10 euros i la B, 16 euros.\begin{tasks}(1)
  \task Com hauré de combinar les dues marques per obtenir la dieta desitjada pel preu mínim?
  \task Quin és el mínim preu que hauré de pagar?
\end{tasks}

\ex{2} \pau{C. Valenciana, juny 2026} Una granja vol crear una mescla fertilitzant amb dos compostos, NutriMax i BioGrow. S'han d'aplicar com a mínim 8 unitats de nitrogen, 6 de fòsfor i 10 de potassi. Un litre de NutriMax conté 2 unitats de nitrogen, 1 de fòsfor i 1 de potassi, i un litre de BioGrow, 1 de nitrogen, 1 de fòsfor i 2 de potassi. El litre de NutriMax costa 5 euros i el de BioGrow, 4 euros.\begin{tasks}(1)
  \task Quants litres de cada producte s'han d'utilitzar perquè el cost de la mescla siga mínim?
  \task Quin és aquest cost mínim?
\end{tasks}


% =====================================================
\seccio{5}{Casos especials: infinites solucions i problemes sense solució}

\ex{2} En la fusteria de l'exercici 16 canvien els preus i ara el benefici és de 30 € per taula i 10 € per cadira. Quina producció dona el màxim benefici? És única?\recorda{Si l'òptim s'assoleix en dos vèrtexs, s'assoleix en tots els punts del segment que els uneix: passa quan la funció objectiu és paral·lela a un costat de la regió.}

\ex{2} En la regió de l'exercici 17 (la dels pinsos), estudia si existeixen el màxim i el mínim de les funcions:\begin{tasks}(2)
  \task $f(x,y)=x+y$
  \task $g(x,y)=2x+y$
\end{tasks}

\ex{2} Una empresa ha de produir $x$ unitats d'un producte A i $y$ d'un producte B. Per capacitat, no pot fabricar més de 10 unitats en total, però un contracte l'obliga a complir $2x+3y\geq40$. Maximitza el benefici $B(x,y)=3x+2y$.

\ex{3} En la regió de l'exercici 6a), troba els valors de $k$ perquè la funció $f(x,y)=kx+y$ assolisca el màxim en infinits punts.

\ex{2} \pau{C. Valenciana, juny 1998} Considera el triangle de vèrtexs $(0,\,0)$, $(2,\,8)$ i $(10,\,3)$. Determina raonadament:\begin{tasks}(1)
  \task El punt del triangle on la funció $f(x,y)=-4x+y+9$ pren el màxim.
  \task El punt del triangle on la funció $g(x,y)=4x+y+12$ pren el màxim.
\end{tasks}

\ex{1} Raona si són vertaderes o falses les afirmacions següents:\begin{tasks}(1)
  \task Si la regió factible és acotada i no buida, la funció objectiu sempre té màxim i mínim.
  \task Si la regió factible no és acotada, la funció objectiu no té mai mínim.
  \task L'òptim d'una funció objectiu no constant pot estar en un punt de l'interior de la regió.
  \task Si dos vèrtexs donen el mateix valor màxim, el problema té infinites solucions.
\end{tasks}

\vfill
\begin{center}\color{muted}\small\faCheckCircle[regular]\ \ Ànim! Comprova els teus resultats amb el resolutor dels apunts interactius de la unitat.\end{center}

\end{document}
